%%%%%%%%%%%%%%%%%%%%%%%%%%%%%%%%%%%%%%%%%%%%%%%%%%%%%%%%%%%%
%% Lecture 10
%%%%%%%%%%%%%%%%%%%%%%%%%%%%%%%%%%%%%%%%%%%%%%%%%%%%%%%%%%%%

\lecture
{10}
{Wednesday, 11 September 2026}
{Introduction to Probability}
{Dr. Nishant Chandgotia}

%%%%%%%%%%%%%%%%%%%%%%%%%%%%%%%%%%%%%%%%%%%%%%%%%%%%%%%%%%%%
%% Hoeffding's Inequality
%%%%%%%%%%%%%%%%%%%%%%%%%%%%%%%%%%%%%%%%%%%%%%%%%%%%%%%%%%%%

\section{Hoeffding's Inequality}

\begin{theorem}{Hoeffding's Inequality}{hoeffding-inequality}

Let $X_1,X_2,\ldots,X_n$ be independent random variables such that
\(
    \mathbb{E}(X_i)=0
    \qquad\text{for all }i,
\)
and
\(
    |X_i|\leq a_i.
\)

Define
\[
    S_n=\sum_{i=1}^{n}X_i.
\]

Then, for every $t>0$,
\[
    \mathbb{P}(S_n\geq t)
    \leq
    \exp\left(
        -\frac{t^2}
        {2\sum_{i=1}^{n}a_i^2}
    \right).
\]

\end{theorem}

\begin{corollary}{Two-Sided Hoeffding Bound}{two-sided-hoeffding}

Under the assumptions of Hoeffding's inequality,
\[
    \mathbb{P}(|S_n|\geq t)
    \leq
    2\exp\left(
        -\frac{t^2}
        {2\sum_{i=1}^{n}a_i^2}
    \right),
    \qquad t>0.
\]

\end{corollary}

\begin{remark}

There are several useful versions of Hoeffding's inequality.

\begin{enumerate}

    \item The inequality can be applied to centered versions of
    random variables. If the random variables do not have mean zero,
    replace $S_n$ by
    \[
        S_n-\sum_{i=1}^{n}\mathbb{E}(X_i).
    \]

    \item There are versions of the inequality for asymmetric
    intervals.

    \item Independence can be relaxed in other concentration
    inequalities. One important example is the
    \emph{Azuma--Hoeffding inequality} for martingales.

\end{enumerate}

\end{remark}

%%%%%%%%%%%%%%%%%%%%%%%%%%%%%%%%%%%%%%%%%%%%%%%%%%%%%%%%%%%%
%% Example: Coin Tosses
%%%%%%%%%%%%%%%%%%%%%%%%%%%%%%%%%%%%%%%%%%%%%%%%%%%%%%%%%%%%

\begin{example}{Coin Tosses}{hoeffding-coin-tosses}

Consider independent coin tosses and let
\[
    X_i=
    \begin{cases}
        1, & \text{if the $i$-th toss is a head},\\
        0, & \text{if the $i$-th toss is a tail}.
    \end{cases}
\]

After centering the variables appropriately, Hoeffding's inequality
gives an exponential bound on deviations of the number of heads from
its expected value.

In particular, when the variables satisfy
\[
    |X_i|\leq 1,
\]
we obtain a bound of the form
\[
    \mathbb{P}(|S_n|\geq t)
    \leq
    2\exp\left(-\frac{t^2}{2n}\right).
\]

For $t=n\alpha$,
\[
    \mathbb{P}(|S_n|\geq n\alpha)
    \leq
    2\exp\left(-\frac{n\alpha^2}{2}\right).
\]

Thus the probability of a deviation of order $n$ decays
exponentially in $n$.

\end{example}

%%%%%%%%%%%%%%%%%%%%%%%%%%%%%%%%%%%%%%%%%%%%%%%%%%%%%%%%%%%%
%% Proof
%%%%%%%%%%%%%%%%%%%%%%%%%%%%%%%%%%%%%%%%%%%%%%%%%%%%%%%%%%%%

\begin{proof}

For $\theta>0$, by Markov's inequality,
\[
\begin{aligned}
    \mathbb{P}(S_n\geq t)
    &=
    \mathbb{P}\left(e^{\theta S_n}\geq e^{\theta t}\right)\\
    &\leq
    \frac{\mathbb{E}\left(e^{\theta S_n}\right)}
    {e^{\theta t}}.
\end{aligned}
\]

Since the $X_i$ are independent,
\[
\mathbb{E}\left(e^{\theta S_n}\right)
=
\prod_{i=1}^{n}
\mathbb{E}\left(e^{\theta X_i}\right).
\]

Because
\(
    -a_i\leq X_i\leq a_i
\)
and $\mathbb{E}(X_i)=0$, convexity of the exponential function gives
\[
    \mathbb{E}\left(e^{\theta X_i}\right)
    \leq
    \frac{1}{2}
    \left(
        e^{-\theta a_i}
        +
        e^{\theta a_i}
    \right).
\]

Using the power-series expansion of the exponential function,
\[
    \frac{e^{-u}+e^u}{2}
    \leq
    e^{u^2/2},
\]
and hence
\[
    \mathbb{E}\left(e^{\theta X_i}\right)
    \leq
    e^{\theta^2a_i^2/2}.
\]

Therefore,
\[
\begin{aligned}
    \mathbb{P}(S_n\geq t)
    &\leq
    \frac{
        \exp\left(
            \dfrac{\theta^2}{2}
            \sum_{i=1}^{n}a_i^2
        \right)
    }
    {e^{\theta t}}\\
    &=
    \exp\left(
        \frac{\theta^2}{2}
        \sum_{i=1}^{n}a_i^2
        -
        \theta t
    \right).
\end{aligned}
\]

Choose
\(
    \theta
    =
    \frac{t}{\sum_{i=1}^{n}a_i^2}.
\)

Then
\[
    \mathbb{P}(S_n\geq t)
    \leq
    \exp\left(
        -\frac{t^2}
        {2\sum_{i=1}^{n}a_i^2}
    \right).
\]

This proves Hoeffding's inequality.

\end{proof}

%%%%%%%%%%%%%%%%%%%%%%%%%%%%%%%%%%%%%%%%%%%%%%%%%%%%%%%%%%%%
%% Independence
%%%%%%%%%%%%%%%%%%%%%%%%%%%%%%%%%%%%%%%%%%%%%%%%%%%%%%%%%%%%

\section{Independence}

Let
\(
    (\Omega,\mathcal{F},\mathbb{P})
\)
be a probability space, and let $I$ be an index set.

\begin{definition}{Independence of Random Variables}{independence-random-variables}

Let
\(
    X_i:(\Omega,\mathcal{F})
    \longrightarrow
    (S,\mathcal{S}),
    \qquad i\in I,
\)
be random variables.

The collection $(X_i)_{i\in I}$ is said to be independent if, for every
finite collection of distinct indices
\(
    i_1,i_2,\ldots,i_n\in I
\)
and measurable sets
\(
    A_1,A_2,\ldots,A_n\in\mathcal{S},
\)
we have
\[
    \mathbb{P}
    \left(
        \bigcap_{j=1}^{n}
        \{X_{i_j}\in A_j\}
    \right)
    =
    \prod_{j=1}^{n}
    \mathbb{P}(X_{i_j}\in A_j).
\]

\end{definition}

%%%%%%%%%%%%%%%%%%%%%%%%%%%%%%%%%%%%%%%%%%%%%%%%%%%%%%%%%%%%
%% Independence of Sigma Algebras
%%%%%%%%%%%%%%%%%%%%%%%%%%%%%%%%%%%%%%%%%%%%%%%%%%%%%%%%%%%%

\begin{definition}{Independence of $\sigma$-Algebras}{independence-sigma-algebras}

Let
\(
    \{\mathcal{F}_i\}_{i\in I}
\)
be a collection of sub-$\sigma$-algebras of $\mathcal{F}$.

They are said to be independent if, for every finite collection of
distinct indices
\(
    i_1,i_2,\ldots,i_n\in I
\)
and sets
\(
    B_{i_j}\in\mathcal{F}_{i_j},
\)
we have
\[
    \mathbb{P}
    \left(
        \bigcap_{j=1}^{n}B_{i_j}
    \right)
    =
    \prod_{j=1}^{n}
    \mathbb{P}(B_{i_j}).
\]

\end{definition}

%%%%%%%%%%%%%%%%%%%%%%%%%%%%%%%%%%%%%%%%%%%%%%%%%%%%%%%%%%%%
%% Sigma Algebra Generated by a Random Variable
%%%%%%%%%%%%%%%%%%%%%%%%%%%%%%%%%%%%%%%%%%%%%%%%%%%%%%%%%%%%

\begin{definition}{$\sigma$-Algebra Generated by a Random Variable}
{sigma-algebra-generated-by-random-variable}

Let
\(
    X:(\Omega,\mathcal{F},\mathbb{P})
    \longrightarrow
    (S,\mathcal{S})
\)
be a random variable.

The $\sigma$-algebra generated by $X$ is
\[
    \sigma(X)
    =
    \sigma
    \left(
        X^{-1}(A):A\in\mathcal{S}
    \right).
\]

Equivalently, $\sigma(X)$ is the smallest $\sigma$-algebra contained
in $\mathcal{F}$ with respect to which $X$ is measurable.

\end{definition}

\begin{observation}

The random variables $(X_i)_{i\in I}$ are independent if and only if
the $\sigma$-algebras
\(
    \bigl(\sigma(X_i)\bigr)_{i\in I}
\)
are independent.

\end{observation}

%%%%%%%%%%%%%%%%%%%%%%%%%%%%%%%%%%%%%%%%%%%%%%%%%%%%%%%%%%%%
%% Elementary Properties
%%%%%%%%%%%%%%%%%%%%%%%%%%%%%%%%%%%%%%%%%%%%%%%%%%%%%%%%%%%%

\begin{proposition}{Basic Properties of Independence}{basic-independence-properties}

Let $B,B'\in\mathcal{F}$.

Then
\[
    B\perp B'
    \quad\Longleftrightarrow\quad
    B\perp (B')^c
    \quad\Longleftrightarrow\quad
    B^c\perp B'.
\]

\end{proposition}

\begin{proposition}{Disjoint Independent Events}{disjoint-independent-events}

Let
\(
    B_1,B_2,\ldots,B_n,B\in\mathcal{F}.
\)

Suppose that the events $B_1,\ldots,B_n$ are pairwise disjoint and
\(
    B_i\perp B
    \qquad\text{for every }i.
\)

Then
\[
    \bigcup_{i=1}^{n}B_i
    \perp B.
\]

\end{proposition}

%%%%%%%%%%%%%%%%%%%%%%%%%%%%%%%%%%%%%%%%%%%%%%%%%%%%%%%%%%%%
%% Question
%%%%%%%%%%%%%%%%%%%%%%%%%%%%%%%%%%%%%%%%%%%%%%%%%%%%%%%%%%%%

\begin{question}{Intersection of Independent Events}{intersection-independent-events}

Suppose
\[
    B_1\perp B
    \qquad\text{and}\qquad
    B_2\perp B.
\]

Does it follow that
\[
    B_1\cap B_2\perp B?
\]

\begin{todo}
Find an explicit counterexample, or use a Venn diagram to determine
whether the statement is true in general.
\end{todo}

\end{question}

%%%%%%%%%%%%%%%%%%%%%%%%%%%%%%%%%%%%%%%%%%%%%%%%%%%%%%%%%%%%
%% Pi-Lambda Systems
%%%%%%%%%%%%%%%%%%%%%%%%%%%%%%%%%%%%%%%%%%%%%%%%%%%%%%%%%%%%

\section{$\pi$-$\lambda$ Systems}

A useful way of checking independence is to reduce the collection of
sets that needs to be checked.

\begin{definition}{$\pi$-System}{pi-system}

Let
\[
    \mathcal{P}\subseteq\mathcal{F}.
\]

We say that $\mathcal{P}$ is a $\pi$-system if it is closed under
finite intersections; that is,
\[
    A,B\in\mathcal{P}
    \quad\Longrightarrow\quad
    A\cap B\in\mathcal{P}.
\]

\end{definition}

\begin{definition}{$\lambda$-System}{lambda-system}

Let
\[
    \mathcal{L}\subseteq\mathcal{F}.
\]

We say that $\mathcal{L}$ is a $\lambda$-system if:

\begin{enumerate}

    \item
    \(
        \Omega\in\mathcal{L}.
    \)

    \item If $A,B\in\mathcal{L}$ and $A\subseteq B$, then
    \(
        B\setminus A\in\mathcal{L}.
    \)

    \item If
    \(
        A_1\subseteq A_2\subseteq\cdots,
        \subseteq 
        A_n, \cdots \in\mathcal{L},
    \)
    then
    \(
        \bigcup_{n=1}^{\infty}A_n\in\mathcal{L}.
    \)

\end{enumerate}

\end{definition}

%%%%%%%%%%%%%%%%%%%%%%%%%%%%%%%%%%%%%%%%%%%%%%%%%%%%%%%%%%%%
%% Pi-Lambda Theorem
%%%%%%%%%%%%%%%%%%%%%%%%%%%%%%%%%%%%%%%%%%%%%%%%%%%%%%%%%%%%

\begin{theorem}{$\pi$-$\lambda$ Theorem}{pi-lambda-theorem}

Let $\mathcal{P}$ be a $\pi$-system and let $\mathcal{L}$ be a
$\lambda$-system such that
\(
    \mathcal{P}\subseteq\mathcal{L}.
\)
Then
\(
    \sigma(\mathcal{P})
    \subseteq
    \mathcal{L}.
\)

\end{theorem}

\begin{proof}
Proof Sketch: Durret
Let $\mathcal{L}'$ be the smallest $\lambda$-system containing
$\mathcal{P}$. Thus,
\[
    \mathcal{P}
    \subseteq
    \mathcal{L}'
    \subseteq
    \mathcal{L}.
\]

The $\pi$-$\lambda$ theorem states that if $\mathcal{P}$ is a
$\pi$-system, then
\[
    \sigma(\mathcal{P})
    =
    \mathcal{L}'.
\]

Therefore,
\[
    \sigma(\mathcal{P})
    \subseteq
    \mathcal{L}.
\]

\end{proof}

%%%%%%%%%%%%%%%%%%%%%%%%%%%%%%%%%%%%%%%%%%%%%%%%%%%%%%%%%%%%
%% Application to Independence
%%%%%%%%%%%%%%%%%%%%%%%%%%%%%%%%%%%%%%%%%%%%%%%%%%%%%%%%%%%%

\begin{remark}

The $\pi$-$\lambda$ theorem allows us to verify independence on a
smaller collection of sets.

Instead of checking independence for every set in
\[
    \sigma(X_i),
\]
it is often enough to check it on a suitable $\pi$-system generating
that $\sigma$-algebra.

\end{remark}

%%%%%%%%%%%%%%%%%%%%%%%%%%%%%%%%%%%%%%%%%%%%%%%%%%%%%%%%%%%%
%% Lecture End
%%%%%%%%%%%%%%%%%%%%%%%%%%%%%%%%%%%%%%%%%%%%%%%%%%%%%%%%%%%%

\lectureend
